\section{Discussion}\label{sec:discussion}

The screen is a first stage. Its intended
pipeline reduces $p$ coordinates to a manageable $d$, after which a joint
method --- the central tail-index subspace of \citet{gardesPodgorny2025}, a
structured tail regression \citep{wangTsai2009,sasakiTaoWang2026}, or
multivariate local estimation --- can resolve interactions and correlated
neighbours among the retained variables. The two stages solve different
problems: rapid elimination of evidently irrelevant coordinates scales to
$p \geq n\alpha$, while accurate joint tail inference does not.

Two limitations remain. First, activity does not imply marginal
detectability. When every fibre of an active coordinate meets the global
argmax set, the score carries no information about that coordinate, at any
sample size. Proposition~\ref{prop:geometry} characterizes this case
exactly. Neither increasing $d$ nor increasing $n$ can overcome this lack
of identification; iterative or grouped extensions would require a
different argument. Second, projection may introduce slowly varying terms with logarithmic
decay, even when the conditional tails are exactly Pareto. Accordingly,
condition (E2) allows the signal to decrease on a logarithmic scale, rather
than at the polynomial rates $\Delta_{\min}\asymp n^{-\kappa}$ commonly
used in SIS theory. This condition is sufficient for the asymptotic result, but does not define a finite-sample detection threshold; the simulations show that the screen can remain effective well before the asymptotic separation regime is reached. The slower signal regime may nevertheless help
explain the comparison with fixed-quantile screening. In A2, where the
tail-index gap is smaller than in A1, Quantile-SIS at a favourable choice
of $\tau$ performs better than the tail-index screen.


In practice, tuning matters more than the theory prescribes. The tail
fraction trades Hill bias against variance; the bandwidth additionally
determines how often near-maximizing configurations of the other active
coordinates enter a window. We recommend the workflow of
Sections~\ref{sec:simulation} and~\ref{sec:realdata}: run the screen on
the block of nine neighbouring tunings and rank by the aggregated
minimum rank. The aggregation is deliberately conservative --- a
predictor is retained when at least one reasonable tuning finds
evidence for it --- which is the right direction of error for a
first-stage sure screen: false exclusion is final, while false
retention is corrected by the second-stage multivariate analysis. The
per-setting ranks describe tuning sensitivity and should accompany the
ranking as a diagnostic, not act as a second screening threshold.
Dependence deserves caution of its own: it can help, by concentrating
active configurations, or hurt, by promoting inactive neighbours, and
the direction is design-specific.

Projection turns a tail-index surface into an upper envelope. Variation
of that envelope is a scalable marginal signal for tail-index activity,
and its absence, characterized by the argmax projections, tells the
analyst exactly what a marginal screen can and cannot see.
